\section{All odd-characteristic symplectic normal cores of rank three}
\label{sec:odd-c3-core}

\paragraph{Status and scope.}
This is a new derivation from the complete official character-table
data and the published automorphism and unipotent-character inputs
specified below. The generic data are from L\"ubeck's dissertation and
its official CHEVIE/GenericCharacterTables implementation; they are
not described as a separately refereed journal table. The new
coefficient and finite-spectrum certificates have actually passed.
The whole mathematical proof below has been reread by its author.
The parent and an independent peer have completed whole-proof,
source and code reviews, all passing; both have executed exact
certificates, and the peer has separately reconstructed all
generic inequalities, finite spectra and auxiliary field identities.
No reality, selfduality,
fixed determinant, or equality of inducing and original character
fields is assumed.

For a finite-order operator $T$ on a complex space of dimension $a$,
put
\[
 \nu(T)=a-\max_{\lambda}\dim\ker(T-\lambda I).
                                                               \tag{C3.1}
\]
This projective support is unchanged by a nonzero scalar multiple.
Ordinary moving rank is at least $\nu(T)$, and a nonidentity scalar
inertia has moving rank $a$. Two elementary inequalities used below
are
\[
 \nu(T)\ge\frac{a-|\operatorname{Tr}T|}{2},\qquad
 \nu(T)\ge\nu(T^j)\quad(j\in\mathbf Z).
                                                               \tag{C3.2}
\]
For the first, rotate an eigenvalue of largest multiplicity to one
and apply the reverse triangle inequality to the remaining
eigenvalues. The second follows by inclusion of each eigenspace of
$T$ in an eigenspace of $T^j$.

\begin{lemma}[D-C3-01: complete similitude data and restriction]
\label{D-C3-01}
Let $q\ge5$ be odd, $C=\operatorname{CSp}_6(q)$, and
$S=\operatorname{Sp}_6(q)$. Every nonlinear ordinary irreducible
character $\chi$ of $C$ satisfies
\[
 \chi(1)\ge q^3-1,\qquad
 |\chi(x)|\le\frac{\chi(1)}4\quad(x\notin Z(C)).
                                                               \tag{C3.3}
\]
Every nontrivial ordinary irreducible character $\rho$ of $S$,
of degree $a$, is a multiplicity-one constituent of an irreducible
$C$-character of degree $a$ or $2a$. Consequently
\[
 a\ge\frac{q^3-1}{2},\qquad
 |\rho(s)|\le a/4
 \quad(s\in S\setminus Z(S)\text{ semisimple}).
                                                               \tag{C3.4}
\]
If $\rho$ is invariant under the diagonal outer automorphism, it
extends to $C$ in the same degree. Every nonscalar operator on its
space implementing an inner-diagonal automorphism then has support
at least $3a/8$.
\end{lemma}

\begin{proof}
The original complete-table input is
\cite[Introduction, pp.~3--4; \S2.2, pp.~18--21;
\S9, pp.~78--79]{C3:Luebeck}, as implemented by the full
\texttt{CSp6.1} table in \cite{C3:GCT}.
Its $100$ character types include every ordinary irreducible
character, and its $90$ class types include every conjugacy class.
The first character type parametrizes the linear similitude
characters. The first class type parametrizes the entire center.
The remaining $99$ character and $89$ class types are retained in
\texttt{generic\_inputs.json}, with their original expressions.
The immutable original source hash is
\texttt{f5e68e9429054370df7376f174a095bc548cf1ecf7ce44f5446f1eaebba250dc}.

Here is the exact finite coefficient certificate for the universal
inequalities (C3.3). Every original table entry is a finite sum
$\sum_j p_j(q)\zeta^{f_j}$, with $p_j\in\mathbf Q[q]$ and each
phase a root of unity at every admissible parameter specialization.
Write $d(q)=\chi(1)$ for the degree polynomial and $q=5+x$. If
$p_j(5+x)=\sum_k c_{jk}x^k$, its absolute value for $x\ge0$ is at
most $\sum_k|c_{jk}|x^k$. For every one of the $8811$ pairs of a
nonlinear row and a noncentral column, the retained rational
coefficient list of
\[
 \frac{d(5+x)}4-\sum_{j,k}|c_{jk}|x^k
                                                               \tag{C3.5}
\]
is nonnegative. Similarly, every nonlinear degree minus $q^3-1$
has nonnegative coefficients after this substitution. Every
noncentral class-length polynomial has nonnegative coefficients
and constant term greater than one. The script
\texttt{check\_generic.py} parses the retained mathematical
expressions without executing the external Julia source and
checks all these coefficients using rational arithmetic. Its
complete lists are in \texttt{generic\_certificate.json}; this
proves (C3.3) for all admissible parameters and all odd $q\ge5$,
not only for finitely many values of $q$.

For the restriction statement, $C/S\simeq\mathbf F_q^\times$
is cyclic, and $C/(SZ(C))$ has order two: the multiplier of a
scalar matrix is its square. The inertia group of $\rho$ therefore
has index one or two in $C$. An irreducible representation of a
normal subgroup extends over a cyclic inertia quotient. Explicitly,
choose an intertwiner for a lift $t$ of its generator; its appropriate
power differs from $\rho(t^k)$ by a scalar, which can be removed by
rescaling. Inducing this extension to $C$ gives an irreducible
character whose restriction is the sum of its distinct conjugates,
each with multiplicity one. Thus its degree is $D=a$ or $2a$.
It is nonlinear because its restriction contains the nontrivial
irreducible $\rho$. The degree assertion in (C3.4) follows.

We verify carefully that semisimple $S$-classes do not split under
$C$-conjugacy. If $s\in S$ is semisimple, the multiplier map
\[
              C_C(s)\longrightarrow\mathbf F_q^\times
                                                               \tag{C3.6}
\]
is surjective. Indeed, use the primary decomposition of its
natural symplectic space. The $+1$ and $-1$ primary spaces are
nondegenerate symplectic spaces and admit a commuting similitude
of every multiplier. A pair of distinct reciprocal primary factors
gives $U\oplus U^*$, on which multiplication by $\mu$ on $U$
and by one on $U^*$ has multiplier $\mu$. For a self-reciprocal
irreducible factor of degree $2d$, the alternating pairing is the
trace of a sesquilinear form over
$\mathbf F_{q^{2d}}/\mathbf F_{q^d}$. Scalar multiplication by
an element of $\mathbf F_{q^{2d}}$ acts on the primary space;
the sesquilinear form is obtained uniquely from the original
pairing by nondegeneracy of the finite-field trace pairing.
The adjoint of this scalar action is the field involution.
Consequently scalar multiplication by
$z$ multiplies this form, and hence the original pairing, by
$z\bar z$. The finite-field norm is surjective, so choose
$z\bar z=\mu$ for any prescribed $\mu\in\mathbf F_q^\times$.
Choosing the same multiplier on all orthogonal summands proves
(C3.6), including primary spaces of any multiplicity. Given a
conjugator in $C$, compose it with a centralizer element of the
inverse multiplier to obtain a conjugator in $S$.

Therefore the one or two constituents of $\chi|_S$ have the same
value at every semisimple $s$, and
$\chi(s)=(D/a)\rho(s)$. This transfers (C3.3) unchanged to the
semisimple assertion in (C3.4). Finally, diagonal invariance makes
the inertia group all of $C$, so the extension has degree $a$.
Any other operator realizing the same automorphism differs from
the extension operator by a scalar, by irreducibility. A nonscalar
extension operator comes from a noncentral natural matrix, so
(C3.2)--(C3.3) give the support $3a/8$.
\end{proof}

\begin{lemma}[D-C3-02: all compatible normalizers]
\label{D-C3-02}
Let $q\ge3$ be odd and let $\rho$ be a nontrivial ordinary
irreducible representation of $\operatorname{Sp}_6(q)$ of degree
$a$. Every nonscalar finite-order linear operator $T$ normalizing
its represented image satisfies the following bounds:
\[
\begin{array}{c|c|c}
q&\text{lower bound for }a&\text{lower bound for }\nu(T)/a\\ \hline
3&13&1/4\\
q\ge5\text{ prime}&(q^3-1)/2&1/5\\
9&364&2/9\\
q\ge25\text{ composite}&(q^3-1)/2&3/16.
\end{array}                                                     \tag{C3.7}
\]
These bounds cover both possible central kernels of $\rho$ and
every compatible diagonal and field automorphism.
\end{lemma}

\begin{proof}
There is no exceptional Schur multiplier for
$\operatorname{PSp}_6(q)$ with odd $q$; see the complete list in
\cite[Chapter~5, Table~1, p.~35]{BI:HHM}.
Its universal perfect cover is $S=\operatorname{Sp}_6(q)$, with
center of order two. Thus every quasisimple representation in
question is already a representation of $S$, with kernel contained
in its center. Automorphisms of either represented central quotient
lift uniquely to $S$. An automorphism of $S$ trivial on its simple
quotient differs from the identity by a homomorphism from the
perfect group $S$ to its center, and hence is the identity.
The standard automorphisms are inner, diagonal, and field, and
\[
 \operatorname{Out}(\operatorname{PSp}_6(q))=C_2\times C_f,
 \qquad q=p^f;
                                                               \tag{C3.8}
\]
the exact source is \cite[\S2.2.2 and Table~2.1,
author p.~14]{HES:LiXia2022}. There is no graph automorphism in
this rank and characteristic.

We also use the published unipotent bound
\[
 \frac{|\rho(u)|}{a}\le\frac1{\sqrt q-1}
 \quad(q>4, u\ne1\text{ unipotent}),
                                                               \tag{C3.9}
\]
with the quasisimple-group conventions and exact radical in
\cite[pp.~2--4, Theorem~2.4(i)]{HES:Liebeck2017}.
This is the published restatement of Gluck's unipotent estimate.

\emph{Prime $q\ge5$.}
There are no field automorphisms. If the diagonal automorphism
stabilizes $\rho$, D-C3-01 gives the stronger support $3a/8$
on the entire compatible normalizer. If it does not, every
compatible normalizer operator is a scalar multiple of an inner
operator. A semisimple inner element has support at least $3a/8$.
An element with a nontrivial unipotent part has a power which is a
nonidentity unipotent of order $p=q$. Let $u$ be that power and
let $\chi$ be the $C$-character in D-C3-01, of degree $D\le2a$.
All $p-1$ nonidentity powers of $u$ are noncentral in $C$.
Fourier inversion and (C3.3) show that every eigenvalue of
$\chi(u)$ has multiplicity at most
\[
 \frac{D+(p-1)D/4}{p}
       =\frac{D(p+3)}{4p}
       \le\frac{a(p+3)}{2p}.
                                                               \tag{C3.10}
\]
Each eigenspace in the constituent $\rho$ is no larger.
Consequently $\nu(\rho(u))\ge a(p-3)/(2p)\ge a/5$.
Taking a power cannot increase the lower bound required for the
original element, by (C3.2). This proves the prime row of (C3.7)
without claiming that restriction is always irreducible.

\emph{Composite $q\ge25$.}
By (C3.9) every nontrivial unipotent has character ratio at most
$1/4$. The semisimple assertion of D-C3-01 and a nontrivial
unipotent power therefore give support at least $3a/8$ for every
noncentral inner element. For a nonscalar normalizer $T$, choose
$h\in S$ for which $[T,\rho(h)]\ne1$; such an $h$ exists by
Schur's lemma. If the commutator is noncentral, its support is
at least $3a/8$, while
\[
 \operatorname{rank}([T,\rho(h)]-I)\le2\nu(T).
                                                               \tag{C3.11}
\]
To check (C3.11), multiply $T$ by a scalar so that
$\operatorname{rank}(T-I)=\nu(T)$ and expand the commutator
difference as the difference of two conjugates of $T-I$.
If the commutator is a nonidentity scalar, its rank is $a$,
which is even stronger. This proves $\nu(T)\ge3a/16$ for
all compatible outer automorphisms, whatever their order.

\emph{$q=9$.}
The same inner argument, now using (C3.9) with value $1/2$,
gives support at least $a/4$ for every noncentral inner operator.
An operator with nontrivial diagonal outer class and trivial field
part, if compatible,
has support at least $3a/8$, by the same-degree extension in
D-C3-01. Inner operators retain the bound $a/4$ just proved.
It remains to handle an automorphism $\alpha$ involving the
field involution $\sigma:x\mapsto x^3$.

If $\alpha^2$ is a nontrivial inner automorphism, $T^2$ is a
scalar multiple of a noncentral inner operator, giving the
stronger bound $a/4$. Suppose instead $\alpha^2=1$. On the
natural space choose a semilinear lift $s=h\sigma$ with
$h\in C=\operatorname{CSp}_6(9)$. Unique lifting to the
universal perfect cover gives
\[
 h\sigma(h)=cI.
                                                               \tag{C3.12}
\]
Applying $\sigma$ and comparing $\sigma(h)h$ shows
$c\in\mathbf F_3^\times$. If $\mu$ is the multiplier of $h$,
then $\operatorname{Nm}_{9/3}(\mu)=c^2=1$.
The norm-one subgroup of $\mathbf F_9^\times$ is precisely
its subgroup of squares. Thus the mixed diagonal-field outer
class cannot have inner-trivial square. This observation is
not needed as a classification assumption in the descent that
follows, but checks that no exceptional mixed involution was
discarded.

Choose $z\in\mathbf F_9^\times$ with
$\operatorname{Nm}_{9/3}(z)=c^{-1}$ and put $s'=zs$.
Then $(s')^2=1$. Its semilinear multiplier is
$\beta=z^2\mu$, of norm one, so choose $d\in\mathbf F_9^\times$
with $d^2=\beta$. If $b$ is the standard symplectic form, then
$b_d=d b$ satisfies
\[
 b_d(s'x,s'y)=\sigma(b_d(x,y)).
                                                               \tag{C3.13}
\]
This follows from $d\beta=d^3$.
Let $\zeta\in\mathbf F_9$ satisfy $\zeta^2=-1$, so
$\sigma(\zeta)=-\zeta$. The fixed space of $s'$ spans the
natural space, since
\[
 x=x_0+\zeta x_1,\qquad
 x_0=\frac{x+s'x}{2},\qquad
 x_1=\frac{x-s'x}{2\zeta},
                                                               \tag{C3.14}
\]
and both $x_0,x_1$ are fixed. Moreover the fixed vectors form a
six-dimensional $\mathbf F_3$-space: a dependence among a fixed
$\mathbf F_3$-basis with coefficients in $\mathbf F_9$ splits
into two $\mathbf F_3$-dependences by (C3.14).
The restriction of $b_d$ is $\mathbf F_3$-valued by (C3.13)
and nondegenerate because this space spans. Choose a symplectic
basis of it. The resulting change of basis $P$ lies in
$\operatorname{CSp}_6(9)$, with multiplier $d^{-1}$, and
$P^{-1}s'P=\sigma$.

Relabel the core representation by this similitude conjugation.
It remains an irreducible representation of the same dimension;
no assertion of its invariance under diagonal conjugation is
needed. The operator $T$ now implements pure $\sigma$. Its
square is scalar, so rescale it to have $T^2=1$.
Inside the standard $\operatorname{GL}_3(9)$ Levi subgroup of
$S$ take
\[
 U=\left\{\begin{pmatrix}1&v\\0&I_2\end{pmatrix}:
              v\in\mathbf F_9^2\right\},\qquad |U|=N=81.
                                                               \tag{C3.15}
\]
The embedded $\operatorname{GL}_2(9)$ acts transitively on the
nonzero additive characters of $U$. Write $t$ for the dimension
of the trivial weight and $m$ for the common multiplicity of
the other $N-1$ weights. Every nonidentity element of $U$ is
unipotent. Averaging (C3.9) gives
\[
 t\le\frac a2+\frac a{2N},\qquad
              m=\frac{a-t}{N-1}\ge\frac a{2N}.
                                                               \tag{C3.16}
\]
Parametrize characters by $w\in\mathbf F_9^2$ using the
$\mathbf F_9/\mathbf F_3$ trace. Frobenius fixes exactly the
$M=9$ parameters in $\mathbf F_3^2$, and pairs the other
$72$. On each paired sum of weight spaces the involution $T$
has $+1$ and $-1$ multiplicities both equal to $m$.
Consequently
\[
 \nu(T)\ge\frac{N-M}{2}\,m
          \ge\frac{N-M}{4N}a=\frac{2a}{9}.
                                                               \tag{C3.17}
\]
The semilinear and dual-character identities are additionally
checked in \texttt{check\_f9.py}; the universal argument is
(C3.12)--(C3.17), not an extrapolation from that finite witness.

\emph{$q=3$.}
The complete retained universal-cover table
\texttt{2.S6(3)} and its cyclic extension
\texttt{2.S6(3).2} are the published ATLAS-origin tables in
\cite{C3:Small}. All $140$ nonlinear core rows and all $88$
extension rows whose restriction is irreducible are included.
Their complete nonscalar spectrum counts and extrema are
\[
\begin{array}{c|r|r|r|c}
 \text{table}&\text{rows}&\text{spectra}&\min a&\min\nu/a\\\hline
 \texttt{2.S6(3)}&140&19460&13&2/7\\
 \texttt{2.S6(3).2}&88&11968&78&8/21.
\end{array}                                                     \tag{C3.18}
\]
The GAP exporter checks full row orthogonality, retains every
original character value and prime power map, and computes all
spectra. A separate exact Python reconstruction selects extension
rows by their restriction inner products and recomputes every
spectrum by rational cyclotomic Fourier arithmetic. It also
checks the fusion, orders, power-map compatibility and complete
normal class unions. Both calculations passed all $31428$
spectra. In particular every nonscalar operator has support at
least $a/4$.

For completeness, the cyclic extension covers every compatible
outer operator, not merely the chosen abstract extension.
An outer automorphism differs from its chosen representative
by an inner automorphism. For that representative take an
intertwiner $A$. Its square differs from the represented
square of a lift in the fixed extension by a scalar, removable
by rescaling $A$. Thus the representation extends to that
fixed extension, and the original operator is scalar times an
element of it. This applies equally when the central involution
acts trivially. There are no other outer automorphisms by
(C3.8). The bounds for $q=3$, together with the other three
cases, exhaust all odd prime powers and prove (C3.7).
\end{proof}

\begin{theorem}[D-C3-03: the complete odd rank-three symplectic gate]
\label{D-C3-03}
Let a finite group $J$ act irreducibly on a complex space $V$.
Suppose it has a normal central product of quasisimple subgroups
acting irreducibly, each with simple quotient
$\operatorname{PSp}_6(q_i)$ for some odd $q_i\ge3$.
The field sizes and factor representations may differ. Every
compatible outer, scalar and factor-permutation action of $J$
is allowed. If this representation is an inducing block of a
finite-group curve-cover construction whose multiplicity
variation has complex dimension six, its period map to a
sixfold Weil domain of dimension nine has rank at most eight.
This holds for every polarization determinant and every
imaginary quadratic original field, even when the inducing
character field is larger.
\end{theorem}

\begin{proof}
Lift each quasisimple factor to its universal cover. Irreducibility
of the central product gives
$V=\bigotimes_i V_i$, with $a_i=\dim V_i\ge13$ and
$e=\dim V=\prod_i a_i$. By D-C3-02 every nonscalar compatible
factor normalizer has support at least $3a_i/16$.
For a factor-preserving normalizer on the tensor product, normalized
eigenvalue multiplicities are convolutions; their maximum cannot
exceed the maximum for any one nonscalar factor. Thus it has
support at least $3e/16$. If all factors are scalar the whole
operator is scalar. A genuine permutation of factors has the
stronger support $e/4$: on a cycle of length $\ell\ge2$ among
equal-dimensional factors, conjugate the tensor intertwiners
so that all twists except the last are identities and diagonalize
the last twist. The weighted rotation of basis words has at most
one vector of each eigenvalue per word orbit. There are $a$
constant words and every other orbit has size at least two,
so the number of orbits is at most
$(a^\ell+a)/2\le3a^\ell/4$. Tensoring with the other cycles
does not increase the normalized largest eigenspace. This also
covers permutations combined with outer actions. Every genuine
nonidentity scalar inertia moves all $e$ dimensions.

Use the actual intermediate finite cover associated to the
inducing block, as in D-BI-01. If its base has genus $h$ and
its faithful-image cover has $R$ genuine branch points, the
six-dimensional multiplicity and the absence of invariant
vectors give
\[
        6=2e(h-1)+D,\qquad
        D=\sum_{j=1}^R\operatorname{rank}(g_j-I).
                                                               \tag{C3.19}
\]
Its variation factors through a finite pointed Hurwitz cover of
$\mathcal M_{h,R}$. Since $e\ge13$, $h\ge2$ is impossible;
if $h=1$, then $D=6$, $R\le6$, and the actual period rank
is at most six. For rank nine one must therefore have
\[
                  h=0,\qquad R\ge12,\qquad D=2e+6.
                                                               \tag{C3.20}
\]
No sufficiency is inferred from this dimension inequality.
Every genuine branch entry has moving rank at least
$\lceil3e/16\rceil$, so
\[
        12\lceil3e/16\rceil\le2e+6,
             \qquad\text{hence}\qquad e\le24.
                                                               \tag{C3.21}
\]
Two or more factors would give $e\ge13^2=169$.
A single factor with $q\ge5$ gives
$e\ge(q^3-1)/2\ge62$. The only remaining case is a single
$q=3$ factor, for which the stronger quarter bound gives
$12\lceil e/4\rceil\le2e+6$, hence $e\le6$, contrary to
$e\ge13$. Thus rank nine is impossible in every stated case.
The argument used one actual complex multiplicity variation and
its intermediate Hurwitz factorization; it never identified an
inducing character field with the original quadratic field.
\end{proof}
